\setcounter{footnote}{49}

\printedpage{39}
% head: A. Grothendieck : Dérivateurs, Ch. XVIII, Édition des Universités de Montpellier II et Paris VII

\begin{itemize}
\item[\textit{1)}] $\Filt'^{\pi}$ \textit{is stable in $\mathcal{M}$ under} $\varinjlim$ \textit{of cardinality $\leqslant \pi$.}
\item[\textit{1$'$)}] $\Filt'^{\pi}$ \textit{is stable under cokernels of double arrows.}
\item[\textit{2)}] \textit{Every object $x$ of $\mathcal{M}$ is of the form $\varinjlim\limits_{I} x_i$, where $x_i$ in} $\Filt'^{\pi}$\textit{, $I$ large relative to $\pi$.}
\item[\textit{3)}] \textit{Condition d) of 1°) implies $P(\pi)$ (hence is equivalent to it), i.e. we have}
\[ \mathcal{M}_{\pi} \ = \ \Filt'^{\pi}. \]
\end{itemize}

\msp{page 127}
\textit{These conditions are satisfied in the following three cases:}
\begin{itemize}
\item[\textit{A)}] \textit{$\pi^{\pi_0} = \pi$ (as a reminder, this is contained in 2°).}
\item[\textit{B)}] \textit{The strict epimorphic families in $\mathcal{M}$ are strict universal epimorphic families.}
\item[\textit{C)}] \textit{Filtering direct limits in $\mathcal{M}$ commute with base change (which is the case in particular if the filtering} $\varinjlim$ \textit{are left exact)} (\footnote{I suppose that this condition amounts to the same condition on $C$, for the filtering $\varinjlim$ of cardinality $\leqslant \pi$.}).
\end{itemize}

\textbf{Remark.} The filtration of $\mathcal{M}$ by the $\Filt'^{\pi}$ is none other than the one considered in SGA 4 I 9.13, where one makes in addition assumption B) on $\mathcal{M}$, which afterwards was no longer so much to my taste. At any rate, according to (4.5.13.1), this filtration comes closer to the ``good'' filtration by the $\mathcal{M}_{\pi}$, than the one I had still been inclined not so long ago to prefer to it, the one by the $\Filt^{\pi}$. It has moreover over the latter the advantage, in the case B) considered, of being a cardinal filtration (which is not necessarily the case for $\Filt^{\pi}$, at least if there exist infinite cardinals $\pi \geqslant \pi_0$ with $\pi^{\pi_0} \neq \pi$). I now suspect that in general (and if there exist infinite cardinals $\pi > \pi_0$ such that $\pi^{\pi_0} \neq \pi$), outside conditions A) B) C), the filtration by the $(\Filt'^{\pi})$ is \textit{not} a
\msp{page 128}
cardinal filtration. In any case, 4.5.6 and the present proposition, 3°, show that if either $(\Filt^{\pi})_{\pi\geqslant\pi_2}$, or $(\Filt'^{\pi})_{\pi\geqslant\pi_2}$ is a cardinal filtration ($\pi_2$ being a given cardinal $\geqslant \pi_1$), then it coincides with the filtration by the $\mathcal{M}_{\pi}$. So apparently, in all the cases where I know how to construct a cardinal filtration, on a paratopos at least (i.e. on a category $\mathcal{M}$ stable under small $\varinjlim$ , and all of whose objects are accessible -- the existence of a cardinal filtration implying moreover that of a small strictly generating subcategory \dots), this is none other than the one by the $\mathcal{M}_{\pi}$ -- something I had not noticed when writing SGA 4 I 9. % [sic] m'étais par aperçu -> m'étais pas aperçu, "par" for the negation "pas"

\textsc{Proof of 4.5.13.}

\textit{Let us prove 1°.} We have the tautological implications
\[
\begin{tikzcd}[row sep=large]
\mathrm{b}'' \arrow[r, Rightarrow] \arrow[d, Rightarrow] & \mathrm{b}' \arrow[r, Rightarrow] \arrow[d, Rightarrow] & \mathrm{b} \arrow[d, Rightarrow] \\
\mathrm{a}'' \arrow[r, Rightarrow] & \mathrm{a}' \arrow[r, Rightarrow] & \mathrm{a} \text{ ,}
\end{tikzcd}
\qquad\qquad \text{(\footnotemark) ,}
\]
\footnotetext{taking into account that $C$ is strictly generating in $\mathcal{M}$, hence $x \simeq \varinjlim\limits_{i\in C/x} x_i$.}
and to prove the equivalence of these six conditions, it suffices to prove
\[ \mathrm{a} \ \Longrightarrow \ \mathrm{b}'' \ . \]
% foot: Version 17 décembre 2025

\printedpage{40}
% head: A. Grothendieck : Dérivateurs, Ch. XVIII, Édition des Universités de Montpellier II et Paris VII

\msp{page 129}
If we have
\[ x \ \simeq \ \varinjlim_{I} x_i \ , \qquad\quad \card \Fl I \ \leqslant \ \pi \ , \]
let $I'$ be the set of finite \textit{subdiagrams} $J \subset I$ of $I$, ordered by inclusion, so $I'$ is a filtering ordered set, and from the fact that $C$ is stable under finite $\varinjlim$, we find a functor
\[ J \longmapsto x_J \ = \ \varinjlim_{i\in J} x_i \ : \ I' \longrightarrow C \ , \]
and we will have
\[ x \ \simeq \ \varinjlim_{J\in I'} x_J \ . \]
On the other hand, we have
\[
\begin{aligned}
\card I' \ &= \ \card \mathfrak{P}_f( \ \underbrace{\Fl(I)}_{\text{of card $\leqslant \pi$}} \ ) \\
&\leqslant \ \pi + \pi^2 + \ \cdots \ + \pi^n + \ \dots \\
&\leqslant \ \pi\aleph_0 \ = \ \pi \ ,
\end{aligned}
\]
q.e.d. To establish 1°, it remains to establish
\[
\begin{tikzcd}[column sep=small]
& P(\pi) \arrow[dl, Rightarrow] \arrow[dr, Rightarrow] & \\
\mathrm{c} & & \mathrm{d} \text{ ,}
\end{tikzcd}
\]
and the reverse implication $\mathrm{c} \Longrightarrow P(\pi)$ when $\pi \geqslant \pi_1$. But $P(\pi) \Longrightarrow \mathrm{c}$ is a tautology, and $P(\pi) \Longrightarrow \mathrm{d}$ since $\mathcal{M}_{\pi}$ is stable under $\varinjlim$ of cardinality $\leqslant \pi$. It remains to prove that $\mathrm{c} \Longrightarrow P(\pi)$ if $\pi \geqslant \pi_1$.

\msp{page 130}
Suppose then that we have a strict epimorphic family
\[ x_{\alpha} \longrightarrow x \qquad\quad (\alpha \in A,\ \card A \leqslant \pi) \ , \]
this therefore means that $x$ is the direct limit of a direct system indexed by the category $B$ having as set of objects $A \amalg (A\times A)$, the nonidentity arrows being represented by the canonical arrows
\[
\begin{tikzcd}[column sep=small]
& x_{\alpha} \times_x x_{\beta} \arrow[dl, "p_{\alpha,\beta}"'] \arrow[dr, "q_{\alpha,\beta}"] & \\
x_{\alpha} & & x_{\beta} \text{ .}
\end{tikzcd}
\]

(\textbf{N.B.} $B$ is not an ordered category, for for the object $(\alpha,\alpha)$, represented by $x_{\alpha}\times_x x_{\alpha}$, % [sic] n'est par une -> n'est pas une, "par" for the negation "pas"
there are \textit{two} arrows into the object $\alpha$, represented by $x_{\alpha}$. But we obviously have $\card \Fl(B) \leqslant \pi$.) If we knew that $C$ is stable under binary products and under passage to subobjects (here the $x_{\alpha} \times_x x_{\beta} \ \subset \ x_{\alpha} \times x_{\beta}$), we would have proved condition a). Failing that, let us again use the fact that $C$ is generating by epimorphisms (even strict ones) and choose for every pair $(\alpha,\beta)$ an epimorphic set of arrows:
\[ \underbrace{x_{\alpha,\beta,\gamma}}_{\in \Ob C_0} \xrightarrow{u_{\alpha,\beta,\gamma}} x_{\alpha} \times_x x_{\beta} \ , \]
% foot: Version 17 décembre 2025

\printedpage{41}
% head: A. Grothendieck : Dérivateurs, Ch. XVIII, Édition des Universités de Montpellier II et Paris VII
where $C_0$ is a reduced exhaustive subcategory of $C$. The arrow $u_{\alpha,\beta,\gamma}$ is known when we know the composites
\msp{page 131}
\[
\begin{tikzcd}
& x_{\alpha} \\
x_{\alpha,\beta,\gamma} \arrow[ur, "v_{\alpha,\beta,\gamma}"] \arrow[dr, "w_{\alpha,\beta,\gamma}"'] & \\
& x_{\beta} \text{ ,}
\end{tikzcd}
\]
so for $x_{\alpha,\beta,\gamma} \in \Ob C_0$ fixed, the possible choices for $u_{\alpha,\beta,\gamma}$ have cardinality bounded by $\pi_1\pi_1 = \pi_1$, moreover we have $\card \Ob C_0 \leqslant \pi_1$, so the total number of arrows $v_{\alpha,\beta,\gamma}$, $w_{\alpha,\beta,\gamma}$ (for $\alpha$, $\beta$ fixed) is bounded by $\pi_1 \cdot \pi_1 = \pi_1$. So the total number of arrows $u_{\alpha,\beta,\gamma}$, for $\alpha$, $\beta$, $\gamma$ variable, is bounded by $\pi \cdot \pi_1 = \pi$ (if $\pi \geqslant \pi_1$). So $x$ is represented by a category $I$ having as objects the disjoint union of $A$ and the set of triple indices $\alpha,\beta,\gamma$, and for each such index, exactly two nonidentity arrows $v_{\alpha,\beta,\gamma}$, $w_{\alpha,\beta,\gamma}$ with this object as source, and going respectively into $x_{\alpha}$, $x_{\beta}$. We thus have $\card \Fl I \leqslant \pi$, and we have a functor $I \longrightarrow C$, $\alpha \longmapsto x_{\alpha}$, $(\alpha,\beta,\gamma) \longmapsto x_{\alpha,\beta,\gamma}$, such that $x \ = \ \varinjlim\limits_{i\in I} x_i$. This completes the proof of 1°).

\msp{page 132}
\textit{Let us prove 2°.} The first inclusion (4.5.13.1) is tautological, the second merely expresses $P(\pi) \Longrightarrow \mathrm{d}$. The fact that $\pi^{\pi_0} = \pi$ implies equality of the extreme terms is a reminder of 4.5.6. It remains to prove that $\Filt'^{\pi}$ is stable in $\mathcal{M}$ under sums of cardinality $\leqslant \pi$. But if $(x_{\alpha})_{\alpha\in A}$ is a family of objects of $\Filt'^{\pi}$, with $\card A \leqslant \pi$, and if $x = \coprod_{\alpha\in A} x_{\alpha}$, taking for every $x_{\alpha}$ a functor $\varphi_{\alpha} : I_{\alpha} \longrightarrow C$, with $\card \Fl I_{\alpha} \leqslant \pi$, such that $x_{\alpha}$ is the $\varinjlim$ in $\mathcal{M}$ of this functor, we find that $x$ is the $\varinjlim$ of the functor $\varphi : I \longrightarrow C$ defined by the $\varphi_{\alpha}$, where $I = \coprod_{\alpha} I_{\alpha}$ (sum in Cat). It is obvious that $\card \Fl(I) \leqslant \pi\cdot\pi = \pi$, which completes the proof of 2°.

\textit{Let us prove 3°.} We have the tautological implications $1 \Longrightarrow 1'$, and $1' \Longrightarrow 1$, since $\Filt'^{\pi}$ is stable under sums of cardinality $\leqslant \pi$. Thus $1 \Longleftrightarrow 1'$, and implies that for every $x$ in $\mathcal{M}$,
\msp{page 133}
$I \ = \ \Filt'^{\pi}/x$ is stable under $\varinjlim$ of cardinality $\leqslant \pi$, hence is large relative to $\pi$. Since $x \simeq \varinjlim\limits_{i\in I} x_i$, this proves that $1 \Longrightarrow 2$. The implication $2 \Longrightarrow 3$ is standard, and follows from $\Filt'^{\pi} \subset \mathcal{M}_{\pi}$, I am going to spell out in a

\textbf{Lemma 4.5.14.} \textit{Let $\mathcal{M}$ be a category stable under small} $\varinjlim$ \textit{, $\pi$ an infinite cardinal, $\mathcal{M}_0$ a strictly full subcategory of $\mathcal{M}$. Consider the conditions}
\begin{itemize}
\item[\textit{a)}] \textit{$\mathcal{M}_0$ is a strictly generating subcategory, and stable in $\mathcal{M}$ under} $\varinjlim$ \textit{of cardinality $\leqslant \pi$.}
\item[\textit{b)}] \textit{For every object of $\mathcal{M}$, we have an isomorphism}
\[ x \ \simeq \ \varinjlim_{i\in I} x_i \ , \]
\textit{the $x_i$ in $\mathcal{M}_0$, $I$ large relative to $\pi$.}
% foot: Version 17 décembre 2025

\printedpage{42}
% head: A. Grothendieck : Dérivateurs, Ch. XVIII, Édition des Universités de Montpellier II et Paris VII
\item[\textit{b$'$)}] \textit{As in b), with $I$ a category that is} ordered.
\item[\textit{c)}] \textit{$\mathcal{M}_{\pi} \subset \mathcal{M}_0$ (where $\mathcal{M}_{\pi}$ is the strictly full subcategory of $\mathcal{M}$ consisting of the $\pi$-accessible objects).}
\end{itemize}
\textit{We thus have the implications}
\[ a \ \Longrightarrow \ b \ \Longleftrightarrow \ b' \ \Longrightarrow \ c \ , \]
\textit{and if we assume $\mathcal{M}_0 \subset \mathcal{M}_{\pi}$, and $\mathcal{M}_{\pi}$ a strictly generating subcategory of $\mathcal{M}$, then the four conditions $a$, $b$, $b'$, $c$ are equivalent, and equivalent to $\mathcal{M}_0 = \mathcal{M}_{\pi}$} (\footnote{cf. proof of proposition 4.5.6, $4 \Longrightarrow 5$, page 117.}).

\msp{page 134}
Finally (returning to the proof of 4.5.13, 3°), the implication $3) \Longrightarrow 1)$ comes from the fact that $\mathcal{M}_{\pi}$ is stable under $\varinjlim$ of cardinality $\leqslant \pi$.

It remains to prove that these conditions, i.e. the stability of $\Filt'^{\pi}$ under cokernels of double arrows, is satisfied in cases A) B) C), case A) being moreover already known. Case B) follows immediately from the transitivity of strict universal epimorphic families, and this is moreover also contained explicitly (stability of $\Filt'^{\pi}$ under $\varinjlim$ of cardinality $\leqslant \pi$) in SGA 4 I 9.13 (if not in the statement, at least it is pointed out in the proof). When writing SGA 4 I 9, I therefore knew that we have $\mathcal{M}_{\pi} \subset \Filt'^{\pi}$ for every $\pi$, but the inclusion, more obvious however in the case at hand, $\Filt'^{\pi} \subset \mathcal{M}_{\pi}$ (when we assume the objects of $C$ $\pi_1$-accessible) had apparently escaped me\,?

It remains to deal with case C). $\Big($I have the impression that it is a little more general than case B), i.e. that condition B) on $\mathcal{M}$ implies C), but not conversely. Thus, condition C) is verified as soon as $\mathcal{M}$ corresponds to a species of algebraic structure defined in
\msp{page 135}
terms of $\varprojlim$ \textit{which are finite}, i.e. when $\mathcal{M}$ is of the form $\Ind(C)$, $C$ stable under $\varinjlim$ \textit{which are finite}, whereas I doubt that condition B) is always satisfied for these categories. For example in the category Cat (which is a paratopos of this species), the filtering $\varinjlim$ are left exact, but I do not believe that the strict epimorphic families are universal. % [sic] épimorphiques stricts -> épimorphiques strictes, agreement with "familles"
Let us take e.g. two group homomorphisms $G_0 \xrightarrow{u_0} G$, $G_1 \xrightarrow{u_1} G$ such that $G_0 * G_1 \longrightarrow G$ is surjective, then this two-element family is strictly epimorphic in Gr and in Cat (\footnote{!, this is \textit{false} in general, but true in the particular case below.}), but is not necessarily universally strictly epimorphic, even if e.g. $G_0 = G_1 = H$, $G = H \times H$, and $u_0, u_1$ are injections $H \rightrightarrows H \times H$, given by $u_0(h) = (h,1)$, and $u_1(h) = (1,h)$. If I take the diagonal map $H \longrightarrow H \times H$, the inverse images of $u_0$, $u_1$ are the same object of $\mathcal{M}/H$, namely $e \longrightarrow H$, and this family is strictly epimorphic only if $H$ is the trivial group.$\Big)$

Case C) will follow from the more precise result:

\textbf{Corollary 4.5.15.} \textit{Under the general conditions of 4.5.13, suppose that in $\mathcal{M}$ the} $\varinjlim$
\msp{page 136}
that are filtering \textit{commute with base change (e.g. the filtering} $\varinjlim$ \textit{are exact). Then for every $\pi \geqslant \pi_1$,} $\Filt'^{\pi}$ \textit{consists of the $x \in \Ob \mathcal{M}$ satisfying the following condition.}
% foot: Version 17 décembre 2025

\printedpage{43}
% head: A. Grothendieck : Dérivateurs, Ch. XVIII, Édition des Universités de Montpellier II et Paris VII
\begin{itemize}
\item[\textit{c$'$)}] \textit{There exists a family of arrows $(u_{\alpha})_{\alpha\in A}$, $u_{\alpha} : x_{\alpha} \longrightarrow x$, in $\mathcal{M}$, which is} universally strictly epimorphic \textit{(i.e. covering for the canonical topology of $\mathcal{M}$), with $x_{\alpha} \in \Ob C$ $\forall\ \alpha \in A$, and} $\card A \leqslant \pi$.
\end{itemize}
\textsc{Proof} : we must prove that this condition is equivalent to condition c) of 4.5.13 1°). Since we trivially have $\mathrm{c}' \Longrightarrow \mathrm{c}$, it remains to show that $\mathrm{c} \Longrightarrow \mathrm{c}'$. But we saw in 4.5.13 1°) that c implies a$''$, i.e. that we have an isomorphism
\[ x \ = \ \varinjlim_{i\in I} x_i \ , \]
the $x_i$ in $C$, and $I$ a \textit{filtering ordered} set, with $\card I \leqslant \pi$. It remains to see that if $x$ in $\mathcal{M}$ is a \textit{filtering} direct limit of objects $x_i$, then the set of arrows $x_i \longrightarrow x$ is strictly epimorphic \textit{universally.} Now this is immediate, for for $x' \longrightarrow x$, we have $x' = \varinjlim\limits_{i} x'_i$ (by assumption C), where $x'_i = x' \times_x x_i$, so the family $x'_i \longrightarrow x'$ is strictly epimorphic, q.e.d.

Let us now prove that in case C, the condition of 3°, 1 is satisfied, i.e. that $\Filt'^{\pi}$ is stable under $\varinjlim$ of cardinality $\leqslant \pi$. This
\msp{page 137}
follows at once from the

\textbf{Corollary 4.5.16.} (of 4.5.15). \textit{Under the conditions of 4.5.15 (i.e. condition C) of 4.5.13, 3°), let $x$ be an object of $\mathcal{M}$, and suppose that there exists a family} epimorphic \underline{universally} $(x_{\alpha} \longrightarrow x)_{\alpha\in A}$ \textit{of arrows of $\mathcal{M}$, with $x_{\alpha} \in$} $\Ob \Filt'^{\pi}$ \textit{for every $\alpha \in A$, and} $\card A \leqslant \pi$. \textit{Then $x \in$} $\Ob \Filt'^{\pi}$.

Indeed, the assumption on $x_{\alpha}$ is equivalent, by virtue of 4.5.15, to the existence of a strict universal epimorphic family of arrows $(x_{\alpha\beta} \longrightarrow x_{\alpha})_{\beta\in B_{\alpha}}$, with $\card B_{\alpha} \leqslant \pi$ and the $x_{\alpha\beta}$ in $C$. Let $B = \coprod_{\alpha\in A} B_{\alpha}$, we then have $\card B \leqslant \pi^2 = \pi$, and on the other hand, the family of composite arrows $x_{\alpha\beta} \longrightarrow x_{\alpha} \longrightarrow x$ indexed by $B$ is universally strictly epimorphic (by transitivity of this type of family), hence $x \in \Ob \Filt'^{\pi}$, q.e.d.

With all that, we have still not given a characterization of $\mathcal{M}_{\pi}$ itself, which is better than $\Filt^{\pi}$ and $\Filt'^{\pi}$ from the categorical point of view, for only the $\mathcal{M}_{\pi}$ define a cardinal filtration, and moreover they have the desirable stability properties in this
\msp{page 138}
context, for $\varinjlim$ of cardinality $\leqslant \pi$. The only nontautological criterion I manage to extract in the general case (for every $\pi \geqslant \pi_0$, and without any additional exactness assumption on $\mathcal{M}$) is the following, much less ``explicit'' or ``computational'' than the very concrete criteria for $\Filt^{\pi}$ and for $\Filt'^{\pi}$:

\textbf{Proposition 4.5.17.} \textit{Let $\pi$ be a cardinal $\geqslant \pi_0$. Then $\mathcal{M}_{\pi}$ is equal to the smallest strictly full subcategory of $\mathcal{M}$, containing $C$, and stable under} $\varinjlim$ \textit{of cardinality $\leqslant \pi$.} (It is important to note that it \textit{is not enough} to say: stable under $\varinjlim\limits_{I}$ , with \textit{$I$ large relative to} $\pi_0$, and $\card I \leqslant \pi$. We then find $\Filt^{\pi}$, which risks being a strictly smaller category than $\mathcal{M}_{\pi}$. Nor is it the category of the $\varinjlim\limits_{I} x_i$, with $x_i \in C$ $\forall\ i \in \Ob I$, and $\card \Fl I \leqslant \pi$ -- since that category is $\Filt'^{\pi}$.)

\textsc{Proof}. Let $\mathcal{M}'$ be the category of the statement, we thus have $\mathcal{M}' \subset \mathcal{M}_{\pi}$, since $\mathcal{M}_{\pi}$ contains $C$ and is stable under the $\varinjlim$ of cardinality $\leqslant \pi$. Conversely, lemma 4.5.14, $\mathrm{a} \Longrightarrow \mathrm{c}$.
% foot: Version 17 décembre 2025